% V8: length calibration — rewritten for compactness (~600 words)
\section{Robustness and Extensions}
\label{sec:robustness}

% V8: length calibration — opening paragraph replaces six subsection intros
Six groups of checks confirm the baseline result. (i)~FL-definition
variants: IQR multiplier, win-rate cutoff, temporal window.
(ii)~Sample and specification variants: unrestricted sample,
homogeneous subsamples, item$\times$year FE, matching.
(iii)~Alternative clustering. (iv)~Sensitivity to unobservables
(Cinelli--Hazlett). (v)~Staggered DiD. (vi)~Welfare bounds.
Table~\ref{tab:robustness_summary} summarises all point estimates.

% V8: length calibration — single summary table replaces six paragraphs of prose
\begin{table}[htbp]
\centering
\caption{Robustness summary: FL price coefficient across checks}
\label{tab:robustness_summary}
\small
\begin{tabular}{lcccc}
\hline\hline
Check & Coeff.\ & SE & $N$ & Note \\
\hline
\multicolumn{5}{l}{\textit{FL definition variants}} \\
IQR 1.0$\times$ & 0.044 & 0.019 & 1{,}673{,}837 & 3{,}442 FL firms \\
IQR 2.0$\times$ & 0.071 & 0.021 & 1{,}673{,}837 & 2{,}235 FL firms \\
IQR 3.0$\times$ & 0.091 & 0.025 & 1{,}673{,}837 & 1{,}456 FL firms \\
Low win-rate ($<1\%$) & 0.062 & 0.021 & 1{,}673{,}837 & Table~\ref{tab:fl_lowwinrate} \\
Temporal FL (3-yr rolling) & 0.070 & 0.022 & 1{,}255{,}998 & Table~\ref{tab:fl_temporal} \\
Cross-fit average & 0.036 & — & — & avg.\ of two folds; Table~\ref{tab:fl_crossfit} \\
\hline
\multicolumn{5}{l}{\textit{Sample and specification}} \\
Unrestricted sample & 0.063 & 0.020 & 2{,}875{,}653 & no FL-item restriction \\
Homogeneous (low CV) & 0.058 & 0.023 & 836{,}918 & below-median price CV \\
Homogeneous (high CV) & 0.069 & 0.027 & 836{,}919 & above-median price CV \\
Item$\times$year FE & 0.074 & 0.022 & 1{,}673{,}837 & tighter controls \\
CEM & 0.077 & 0.024 & — & primary result (Sec.~\ref{sec:results_ols}) \\
IPW & 0.055 & 0.021 & — & primary result (Sec.~\ref{sec:results_ols}) \\
\hline
\multicolumn{5}{l}{\textit{Alternative clustering}} \\
Item clustering (baseline) & 0.064 & 0.022 & 1{,}673{,}837 & \\
PBU clustering & 0.064 & 0.014 & 1{,}673{,}837 & \\
Two-way (item + PBU) & 0.064 & 0.024 & 1{,}673{,}837 & \\
\hline\hline
\end{tabular}
\end{table}

% V8: length calibration — Cinelli-Hazlett compressed to one paragraph
\paragraph{Sensitivity to unobservables.}
\citet{cinelli2020making} sensitivity analysis yields a robustness
value $RV_{q=1} = 17.5\%$: an unobserved confounder would need to
explain at least 17.5\% of residual variation in \emph{both} FL
presence and log prices to nullify the coefficient---a stringent
benchmark given the rich item, year, and PBU fixed-effects structure.
Full sensitivity plots are in Appendix~\ref{app:robustness}.

% V8: length calibration — staggered DiD compressed to one paragraph
\paragraph{Staggered difference-in-differences.}
We report the \citet{callaway2021difference} staggered DiD as a
complementary, underpowered exercise. The panel comprises 144,168
market-year observations (1,511 treated, 18,266 never-treated
markets). The overall ATT is 0.014 (SE~=~0.039), positive but
insignificant. The minimum detectable effect is
$1.96\times 0.039 = 0.076$, which exceeds our OLS estimate of 0.064;
the design therefore lacks power to detect effects of the magnitude
estimated by OLS and IV. \citet{rambachan2023more} sensitivity
analysis, sharp-entry subsamples, and event-study plots are in
Appendix~\ref{app:did} (Table~\ref{tab:did_revised},
Figure~\ref{fig:event_study}).

% V8: length calibration — welfare paragraph shortened; detail stays in Sec 7.3
\subsection{Welfare Bounds}
\label{sec:welfare}

Welfare loss is computed as $(\exp(\hat{\beta})-1)\times\sum_{i:\,\text{FL}_i=1}p_i$.
The OLS-based lower bound is R\$734M (95\%~CI: R\$243M--R\$1,246M),
representing 2.4\% of total BEC procurement over 2009--2019. The
cross-fit average (0.036) yields a conservative intermediate estimate
of R\$400M (1.3\% of spending), breaking any mechanical link between
FL classification and outcomes. The IV-based upper bound is
R\$2.4B (7.8\%), reflecting the instrumented LATE for complier
tenders. Full derivation and bounds table are in
Table~\ref{tab:welfare_bounds}, Appendix~\ref{app:extensions}.

% V8: length calibration — oversight section retained as in v7 (already compact)
\subsection{Oversight Heterogeneity} % V7-T2: honest power discussion

Proposition~\ref{prop:optimal_m} predicts that $m^*$ is decreasing in
detection probability $\theta_k$, implying a larger FL--price effect
at PBUs with weaker oversight. Table~\ref{tab:regime_oversight} in
Appendix~\ref{app:robustness} reports the FL price coefficient by PBU
size quartile (proxying oversight intensity with annual procurement
volume). The point estimates decrease monotonically from Q1 (0.214)
to Q4 (0.017), consistent with the prediction, but only Q3 (0.066)
is individually significant ($p < 0.05$). We interpret this pattern
as \emph{suggestive but underpowered}: the wide confidence intervals
for Q1 and Q2 preclude strong inference from the quartile split. The
counterfactual analysis (Section~\ref{sec:counterfactual}, Panel~C)
provides a complementary test using the semi-elasticity of FL count
with respect to PBU size as a continuous proxy for
$\partial m^*/\partial\theta$.

\begin{figure}[htbp]
\centering
\includegraphics[width=0.85\textwidth]{fig_07_threshold_stability.pdf}
\caption{FL price coefficient across IQR multiplier thresholds.
Error bars show 95\% confidence intervals.}
\label{fig:threshold_stability}
\end{figure}

The IQR classification rule is compared with model-based alternatives % V8: brief ref, detail in appendix
(structural likelihood ratio, random forest with 7 bid-level features)
in Appendix~\ref{app:robustness}. The IQR rule (AUC~$= 0.94$)
dominates both alternatives (DeLong $p = 0.04$ vs.\ random forest
AUC~$= 0.90$), validating the threshold rule as the most discriminative
classifier available.
